% ch2-applications-lineaires.tex — Chapitre 2 : Applications linéaires et matrices
% Ce chapitre introduit la notion centrale d'application linéaire, en
% développe les propriétés fondamentales (noyau, image, théorème du rang),
% puis établit le lien avec les matrices et les changements de base.

\chapter{Applications linéaires et matrices}\label{ch:applications-lineaires}

\section*{Introduction}
\addcontentsline{toc}{section}{Introduction}

L'algèbre linéaire est, en un sens, l'étude des \emph{applications
linéaires}\index{application linéaire}~: ces transformations qui respectent
la structure d'espace vectoriel. Avant d'en donner la définition formelle,
observons quelques exemples géométriques dans le plan~$\R^2$.

\medskip

Considérons les transformations suivantes du plan~:
\begin{itemize}
  \item la \textbf{rotation}\index{rotation} d'angle $\theta$ autour de
    l'origine, qui envoie le vecteur $\begin{pmatrix}x\\y\end{pmatrix}$
    sur $\begin{pmatrix}x\cos\theta-y\sin\theta\\x\sin\theta+y\cos\theta\end{pmatrix}$~;
  \item la \textbf{réflexion}\index{réflexion} par rapport à l'axe des
    abscisses, qui envoie $\begin{pmatrix}x\\y\end{pmatrix}$ sur
    $\begin{pmatrix}x\\-y\end{pmatrix}$~;
  \item la \textbf{projection}\index{projection} orthogonale sur l'axe des
    abscisses, qui envoie $\begin{pmatrix}x\\y\end{pmatrix}$ sur
    $\begin{pmatrix}x\\0\end{pmatrix}$~;
  \item l'\textbf{homothétie}\index{homothétie} de rapport $\lambda$, qui
    envoie $\begin{pmatrix}x\\y\end{pmatrix}$ sur
    $\begin{pmatrix}\lambda x\\\lambda y\end{pmatrix}$.
\end{itemize}

\medskip

Ces quatre transformations partagent deux propriétés remarquables~: elles
envoient la somme de deux vecteurs sur la somme des images, et le produit
d'un vecteur par un scalaire sur le produit du scalaire par l'image. C'est
précisément cette double compatibilité qui caractérise les applications
linéaires.

\begin{figure}[ht]
\centering
\begin{tikzpicture}[scale=1.2, >=Stealth]
  % Grille originale
  \begin{scope}
    \draw[gray!30, very thin] (-0.5,-0.5) grid (2.5,2.5);
    \draw[->] (-0.5,0) -- (2.7,0) node[right] {$x$};
    \draw[->] (0,-0.5) -- (0,2.7) node[above] {$y$};
    \draw[->, thick, blue] (0,0) -- (1,0) node[below right] {$\vece{1}$};
    \draw[->, thick, red] (0,0) -- (0,1) node[above left] {$\vece{2}$};
    \draw[->, thick, black!60!green] (0,0) -- (1,1) node[above right] {$\vec{v}$};
    \node[below] at (1.2,-0.7) {Vecteurs originaux};
  \end{scope}
  % Rotation
  \begin{scope}[xshift=5cm]
    \draw[gray!30, very thin] (-1.5,-0.5) grid (2.5,2.5);
    \draw[->] (-1.5,0) -- (2.7,0) node[right] {$x$};
    \draw[->] (0,-0.5) -- (0,2.7) node[above] {$y$};
    \draw[->, thick, blue] (0,0) -- ({cos(45)},{sin(45)}) node[above right] {$f(\vece{1})$};
    \draw[->, thick, red] (0,0) -- ({-sin(45)},{cos(45)}) node[above left] {$f(\vece{2})$};
    \draw[->, thick, black!60!green] (0,0) -- (0,{sqrt(2)}) node[above right] {$f(\vec{v})$};
    \node[below] at (1.0,-0.7) {Rotation de $\pi/4$};
  \end{scope}
\end{tikzpicture}
\caption{Effet d'une rotation de $\pi/4$ sur les vecteurs du plan.}
\label{fig:rotation-plan}
\end{figure}

\begin{figure}[ht]
\centering
\begin{tikzpicture}[scale=1.2, >=Stealth]
  % Projection
  \begin{scope}
    \draw[gray!30, very thin] (-0.5,-1.5) grid (2.5,2.5);
    \draw[->] (-0.5,0) -- (2.7,0) node[right] {$x$};
    \draw[->] (0,-1.5) -- (0,2.7) node[above] {$y$};
    \draw[->, thick, blue] (0,0) -- (2,1.5) node[above right] {$\vec{v}$};
    \draw[->, thick, red] (0,0) -- (2,0) node[below right] {$p(\vec{v})$};
    \draw[dashed, gray] (2,1.5) -- (2,0);
    \node[below] at (1.2,-1.8) {Projection sur $Ox$};
  \end{scope}
  % Cisaillement
  \begin{scope}[xshift=6cm]
    \draw[gray!30, very thin] (-0.5,-1.5) grid (3.5,2.5);
    \draw[->] (-0.5,0) -- (3.7,0) node[right] {$x$};
    \draw[->] (0,-1.5) -- (0,2.7) node[above] {$y$};
    % Carré original
    \draw[blue, thick] (0,0) -- (1,0) -- (1,1) -- (0,1) -- cycle;
    % Cisaillement k=1
    \draw[red, thick] (0,0) -- (1,0) -- (2,1) -- (1,1) -- cycle;
    \draw[->, thick, blue] (0,0) -- (1,0) node[below] {$\vece{1}$};
    \draw[->, thick, blue] (0,0) -- (0,1) node[left] {$\vece{2}$};
    \draw[->, thick, red] (0,0) -- (1,0);
    \draw[->, thick, red] (0,0) -- (1,1) node[above right] {$f(\vece{2})$};
    \node[below] at (1.5,-1.8) {Cisaillement ($k=1$)};
  \end{scope}
\end{tikzpicture}
\caption{Projection orthogonale et cisaillement~: deux transformations linéaires.}
\label{fig:projection-cisaillement}
\end{figure}

% ============================================================
\section{Définition d'une application linéaire}\label{sec:def-app-lin}
% ============================================================

\begin{definition}{Application linéaire}{def:app-lin}
  Soient $E$ et $F$ deux $\K$-espaces vectoriels. Une application
  $f\colon E\to F$ est dite \emph{linéaire}\index{application linéaire}
  si elle vérifie les deux conditions suivantes~:
  \begin{enumerate}[label=(\roman*)]
    \item \textbf{Additivité~:} $\forall\, \vec{u},\vec{v}\in E,\quad
      f(\vec{u}+\vec{v})=f(\vec{u})+f(\vec{v})$.
    \item \textbf{Homogénéité~:} $\forall\,\lambda\in\K,\;\forall\,
      \vec{u}\in E,\quad f(\lambda\vec{u})=\lambda\, f(\vec{u})$.
  \end{enumerate}
  De manière équivalente, $f$ est linéaire si et seulement si~:
  \[
    \forall\,\lambda,\mu\in\K,\;\forall\,\vec{u},\vec{v}\in E,\quad
    f(\lambda\vec{u}+\mu\vec{v})=\lambda\,f(\vec{u})+\mu\,f(\vec{v}).
  \]
  On note $\cL(E,F)$\index{L(E,F)@$\cL(E,F)$} l'ensemble des applications
  linéaires de $E$ dans $F$. Lorsque $E=F$, on parle d'\emph{endomorphisme}%
  \index{endomorphisme} et l'on note $\End(E)\deff\cL(E,E)$.
\end{definition}

\begin{remark}{Terminologie}{rem:terminologie-al}
  On emploie parfois les termes \emph{morphisme d'espaces vectoriels},
  \emph{opérateur linéaire} (surtout lorsque $E=F$) ou
  \emph{homomorphisme}\index{homomorphisme}. En anglais, on dit
  \emph{linear map} ou \emph{linear transformation}.
\end{remark}

\begin{proposition}{Propriétés élémentaires}{prop:elem-al}
  Soit $f\colon E\to F$ une application linéaire. Alors~:
  \begin{enumerate}[label=(\roman*)]
    \item $f(\vzero_E)=\vzero_F$.
    \item $\forall\,\vec{u}\in E,\quad f(-\vec{u})=-f(\vec{u})$.
    \item Pour toute combinaison linéaire finie~:
      $f\!\left(\displaystyle\sum_{i=1}^{n}\lambda_i\vec{u}_i\right)
      =\sum_{i=1}^{n}\lambda_i\,f(\vec{u}_i)$.
  \end{enumerate}
\end{proposition}

\begin{proof}
  \emph{(i)} On a $f(\vzero_E)=f(0\cdot\vzero_E)=0\cdot f(\vzero_E)
  =\vzero_F$.

  \emph{(ii)} $f(-\vec{u})=f((-1)\cdot\vec{u})=(-1)\cdot f(\vec{u})
  =-f(\vec{u})$.

  \emph{(iii)} Par récurrence immédiate à partir de l'additivité et de
  l'homogénéité.
\end{proof}

\begin{remark}{Détermination par les images d'une base}{rem:determination-base}
  Si $E$ est de dimension finie et si $\cB=(\vece{1},\ldots,\vece{n})$
  est une base de~$E$, alors une application linéaire $f\colon E\to F$
  est entièrement déterminée par les images $f(\vece{1}),\ldots,f(\vece{n})$.
  En effet, pour $\vec{u}=\sum_{i=1}^{n}x_i\,\vece{i}$, on a
  $f(\vec{u})=\sum_{i=1}^{n}x_i\,f(\vece{i})$. Réciproquement, la
  donnée de $n$ vecteurs quelconques $\vecf{1},\ldots,\vecf{n}\in F$
  définit une unique application linéaire $f\in\cL(E,F)$ telle que
  $f(\vece{i})=\vecf{i}$ pour tout~$i$.
\end{remark}

% ============================================================
\section{Exemples fondamentaux}\label{sec:exemples-al}
% ============================================================

\begin{example}{Rotation dans $\R^2$}{ex:rotation}
  Soit $R_\theta\colon\R^2\to\R^2$ la rotation d'angle~$\theta$ autour de
  l'origine. Pour $\vec{u}=\begin{pmatrix}x\\y\end{pmatrix}$, on a
  \[
    R_\theta(\vec{u})=\begin{pmatrix}
      x\cos\theta-y\sin\theta \\ x\sin\theta+y\cos\theta
    \end{pmatrix}.
  \]
  Vérifions la linéarité. Soient $\vec{u}=\begin{pmatrix}x_1\\y_1\end{pmatrix}$,
  $\vec{v}=\begin{pmatrix}x_2\\y_2\end{pmatrix}$ et $\lambda\in\R$. Alors~:
  \begin{align*}
    R_\theta(\vec{u}+\vec{v})
    &=\begin{pmatrix}(x_1+x_2)\cos\theta-(y_1+y_2)\sin\theta\\
      (x_1+x_2)\sin\theta+(y_1+y_2)\cos\theta\end{pmatrix}
    =R_\theta(\vec{u})+R_\theta(\vec{v}),\\
    R_\theta(\lambda\vec{u})
    &=\begin{pmatrix}\lambda x_1\cos\theta-\lambda y_1\sin\theta\\
      \lambda x_1\sin\theta+\lambda y_1\cos\theta\end{pmatrix}
    =\lambda\, R_\theta(\vec{u}).
  \end{align*}
  Donc $R_\theta$ est un endomorphisme de~$\R^2$.
\end{example}

\begin{example}{Projection orthogonale}{ex:projection}
  La projection orthogonale $p\colon\R^2\to\R^2$ sur l'axe des abscisses
  est définie par $p\begin{pmatrix}x\\y\end{pmatrix}
  =\begin{pmatrix}x\\0\end{pmatrix}$. C'est une application linéaire~:
  on vérifie aisément que $p(\vec{u}+\vec{v})=p(\vec{u})+p(\vec{v})$ et
  $p(\lambda\vec{u})=\lambda\,p(\vec{u})$. De plus, $p\circ p=p$~: c'est
  un \emph{projecteur}\index{projecteur} (ou \emph{idempotent}).
\end{example}

\begin{example}{Dérivation}{ex:derivation}
  Soit $E=\R_n[X]$ l'espace des polynômes de degré au plus~$n$ à
  coefficients réels. L'application \emph{dérivation}\index{dérivation}
  \[
    D\colon\R_n[X]\to\R_{n-1}[X],\qquad P\mapsto P'
  \]
  est linéaire~: $(P+Q)'=P'+Q'$ et $(\lambda P)'=\lambda P'$. Son noyau
  est l'ensemble des polynômes constants et son image est $\R_{n-1}[X]$
  tout entier.
\end{example}

\begin{example}{Intégration}{ex:integration}
  L'application
  \[
    I\colon\mathcal{C}([0,1],\R)\to\R,\qquad f\mapsto\int_0^1 f(t)\,dt
  \]
  est linéaire (par linéarité de l'intégrale de Riemann). Ici, l'espace
  de départ est de dimension infinie et l'espace d'arrivée est de
  dimension~$1$.
\end{example}

\begin{example}{Multiplication par une matrice}{ex:mult-matrice}
  Soit $A\in\Mat_{m,n}(\K)$ une matrice à $m$ lignes et $n$ colonnes.
  L'application
  \[
    f_A\colon\K^n\to\K^m,\qquad \vec{x}\mapsto A\vec{x}
  \]
  est linéaire~: $A(\vec{x}+\vec{y})=A\vec{x}+A\vec{y}$ et
  $A(\lambda\vec{x})=\lambda\,A\vec{x}$. Nous verrons que,
  réciproquement, toute application linéaire entre espaces de dimension
  finie peut se représenter par une telle multiplication matricielle
  (
ef{sec:matrice-al}).
\end{example}

\begin{example}{Trace d'une matrice}{ex:trace-lin}
  L'application
  \[
    \tr\colon\Mat_n(\K)\to\K,\qquad A=(a_{ij})\mapsto\sum_{i=1}^n a_{ii}
  \]
  est linéaire~: $\tr(A+B)=\tr(A)+\tr(B)$ et
  $\tr(\lambda A)=\lambda\,\tr(A)$.
\end{example}

% ============================================================
\section{Noyau et image}\label{sec:noyau-image}
% ============================================================

\begin{definition}{Noyau}{def:noyau}
  Soit $f\colon E\to F$ une application linéaire. Le \emph{noyau}%
  \index{noyau} de $f$ est
  \[
    \Ker f\deff\setcond{\vec{u}\in E}{f(\vec{u})=\vzero_F}.
  \]
\end{definition}

\begin{definition}{Image}{def:image-al}
  Soit $f\colon E\to F$ une application linéaire. L'\emph{image}%
  \index{image} de $f$ est
  \[
    \im f\deff\setcond{f(\vec{u})}{\vec{u}\in E}
    =\setcond{\vec{w}\in F}{\exists\,\vec{u}\in E,\;f(\vec{u})=\vec{w}}.
  \]
\end{definition}

\begin{theorem}{Le noyau et l'image sont des sous-espaces vectoriels}{thm:ker-im-sev}
  Soit $f\colon E\to F$ une application linéaire.
  \begin{enumerate}[label=(\roman*)]
    \item $\Ker f$ est un sous-espace vectoriel de~$E$.
    \item $\im f$ est un sous-espace vectoriel de~$F$.
  \end{enumerate}
\end{theorem}

\begin{proof}
  \emph{(i)} Le noyau est non vide car $f(\vzero_E)=\vzero_F$, donc
  $\vzero_E\in\Ker f$. Soient $\vec{u},\vec{v}\in\Ker f$ et
  $\lambda\in\K$. Alors
  \[
    f(\vec{u}+\lambda\vec{v})=f(\vec{u})+\lambda\,f(\vec{v})
    =\vzero_F+\lambda\,\vzero_F=\vzero_F,
  \]
  donc $\vec{u}+\lambda\vec{v}\in\Ker f$. Ainsi $\Ker f$ est un
  sous-espace vectoriel de~$E$.

  \emph{(ii)} L'image est non vide car $\vzero_F=f(\vzero_E)\in\im f$.
  Soient $\vec{w}_1,\vec{w}_2\in\im f$ et $\lambda\in\K$. Il existe
  $\vec{u}_1,\vec{u}_2\in E$ tels que $f(\vec{u}_1)=\vec{w}_1$ et
  $f(\vec{u}_2)=\vec{w}_2$. Alors
  \[
    \vec{w}_1+\lambda\vec{w}_2=f(\vec{u}_1)+\lambda\,f(\vec{u}_2)
    =f(\vec{u}_1+\lambda\vec{u}_2)\in\im f.
  \]
  Donc $\im f$ est un sous-espace vectoriel de~$F$.
\end{proof}

\begin{figure}[ht]
\centering
\begin{tikzpicture}[>=Stealth, scale=0.9]
  % E
  \draw[thick, blue!60, fill=blue!5, rounded corners=8pt]
    (-3,-2.5) rectangle (1.5,2.5);
  \node[blue!80, font=\large\bfseries] at (-0.75,2.8) {$E$};
  % Ker f
  \draw[thick, red!70, fill=red!10, rounded corners=5pt]
    (-2.5,-1.5) rectangle (-0.3,1.5);
  \node[red!80, font=\bfseries] at (-1.4,1.8) {$\Ker f$};
  \node[red!60] at (-1.4,0) {$f=\vzero$};
  % F
  \draw[thick, blue!60, fill=blue!5, rounded corners=8pt]
    (4,-2.5) rectangle (8.5,2.5);
  \node[blue!80, font=\large\bfseries] at (6.25,2.8) {$F$};
  % Im f
  \draw[thick, black!60!green, fill=green!8, rounded corners=5pt]
    (4.8,-1.5) rectangle (7.7,1.5);
  \node[black!60!green, font=\bfseries] at (6.25,1.8) {$\im f$};
  % Flèche
  \draw[->, very thick, black!70] (1.8,0.5) to[bend left=15]
    node[above, font=\small] {$f$} (4.5,0.5);
  \draw[->, very thick, black!70] (1.8,-0.5) to[bend right=15]
    node[below, font=\small] {} (4.5,-0.5);
  % Points
  \fill[red!70] (-1.4,-0.5) circle (2pt) node[left, font=\small] {$\vec{u}$};
  \fill[black!60!green] (6.25,0) circle (2pt) node[right, font=\small] {$\vzero_F$};
  \fill[blue!70] (0.5,0.8) circle (2pt) node[right, font=\small] {$\vec{v}$};
  \fill[black!60!green] (6.0,0.7) circle (2pt) node[right, font=\small] {$f(\vec{v})$};
\end{tikzpicture}
\caption{Noyau et image d'une application linéaire $f\colon E\to F$.}
\label{fig:noyau-image}
\end{figure}

\begin{definition}{Rang}{def:rang}
  Soit $f\colon E\to F$ une application linéaire. Le \emph{rang}%
  \index{rang} de $f$ est la dimension de son image~:
  \[
    \rg(f)\deff\dim(\im f).
  \]
\end{definition}

\begin{proposition}{Image d'une famille génératrice}{prop:image-generatrice}
  Soit $f\colon E\to F$ linéaire. Si $(u_1,\ldots,u_p)$ engendre~$E$,
  alors $(f(u_1),\ldots,f(u_p))$ engendre $\im f$. En particulier, si
  $\cB=(e_1,\ldots,e_n)$ est une base de~$E$, alors
  \[
    \im f=\Vect(f(e_1),\ldots,f(e_n)).
  \]
\end{proposition}

\begin{proof}
  Soit $\vec{w}\in\im f$. Il existe $\vec{u}\in E$ tel que
  $f(\vec{u})=\vec{w}$. Comme $(u_1,\ldots,u_p)$ engendre~$E$, on peut
  écrire $\vec{u}=\sum_{i=1}^p\lambda_i u_i$, d'où
  $\vec{w}=f(\vec{u})=\sum_{i=1}^p\lambda_i f(u_i)$.
\end{proof}

% ============================================================
\section{Théorème du rang}\label{sec:thm-rang}
% ============================================================

Le résultat suivant est l'un des théorèmes les plus importants de
l'algèbre linéaire. Il relie la dimension de l'espace de départ, celle
du noyau et celle de l'image.

\begin{theorem}{Théorème du rang}{thm:rang}
  Soit $f\colon E\to F$ une application linéaire avec $E$ de dimension
  finie. Alors $\Ker f$ et $\im f$ sont de dimension finie et
  \[
    \boxed{\dim E = \dim\Ker f + \dim\im f = \dim\Ker f + \rg(f).}
  \]
\end{theorem}
\index{théorème du rang}

\begin{proof}
  Posons $n=\dim E$ et $r=\dim\Ker f$.

  \medskip
  \textbf{Étape 1.} Le noyau $\Ker f$ est un sous-espace vectoriel de~$E$
  (
ef{thm:thm:ker-im-sev}), donc il est de dimension finie avec $r\leq n$.
  Choisissons une base $(\vec{u}_1,\ldots,\vec{u}_r)$ de~$\Ker f$.

  \medskip
  \textbf{Étape 2.} D'après le théorème de la base incomplète, on peut
  compléter cette base en une base de~$E$~:
  \[
    \cB=(\vec{u}_1,\ldots,\vec{u}_r,\vec{v}_1,\ldots,\vec{v}_s)
  \]
  où $r+s=n$, c'est-à-dire $s=n-r$.

  \medskip
  \textbf{Étape 3.} Montrons que $(f(\vec{v}_1),\ldots,f(\vec{v}_s))$ est
  une base de $\im f$. Cela établira $\dim\im f=s=n-r$, d'où le résultat.

  \medskip
  \textbf{Étape 3a --- Famille génératrice.} Soit $\vec{w}\in\im f$. Il
  existe $\vec{x}\in E$ tel que $f(\vec{x})=\vec{w}$. Décomposons
  $\vec{x}$ dans la base~$\cB$~:
  \[
    \vec{x}=\sum_{i=1}^r\alpha_i\,\vec{u}_i+\sum_{j=1}^s\beta_j\,\vec{v}_j.
  \]
  Par linéarité de~$f$~:
  \[
    \vec{w}=f(\vec{x})
    =\sum_{i=1}^r\alpha_i\,\underbrace{f(\vec{u}_i)}_{=\,\vzero_F}
     +\sum_{j=1}^s\beta_j\,f(\vec{v}_j)
    =\sum_{j=1}^s\beta_j\,f(\vec{v}_j).
  \]
  Donc $(f(\vec{v}_1),\ldots,f(\vec{v}_s))$ engendre $\im f$.

  \medskip
  \textbf{Étape 3b --- Liberté.} Supposons que
  $\sum_{j=1}^s\beta_j\,f(\vec{v}_j)=\vzero_F$. Par linéarité, cela
  signifie $f\!\left(\sum_{j=1}^s\beta_j\,\vec{v}_j\right)=\vzero_F$,
  c'est-à-dire $\sum_{j=1}^s\beta_j\,\vec{v}_j\in\Ker f$. Il existe
  donc des scalaires $\alpha_1,\ldots,\alpha_r\in\K$ tels que
  \[
    \sum_{j=1}^s\beta_j\,\vec{v}_j
    =\sum_{i=1}^r\alpha_i\,\vec{u}_i,
  \]
  soit
  \[
    \sum_{i=1}^r(-\alpha_i)\,\vec{u}_i+\sum_{j=1}^s\beta_j\,\vec{v}_j
    =\vzero_E.
  \]
  Comme $(\vec{u}_1,\ldots,\vec{u}_r,\vec{v}_1,\ldots,\vec{v}_s)$ est
  une base de~$E$, tous les coefficients sont nuls~: en particulier,
  $\beta_1=\cdots=\beta_s=0$. La famille $(f(\vec{v}_1),\ldots,f(\vec{v}_s))$
  est donc libre.

  \medskip
  \textbf{Conclusion.} $(f(\vec{v}_1),\ldots,f(\vec{v}_s))$ est une base
  de $\im f$, d'où $\dim\im f=s=n-r=\dim E-\dim\Ker f$.
\end{proof}

\begin{corollary}{Formule du rang}{cor:formule-rang}
  Si $\dim E=n$ et $\dim F=m$, alors pour toute application linéaire
  $f\colon E\to F$~:
  \[
    \rg(f)\leq\min(n,m).
  \]
\end{corollary}

\begin{proof}
  On a $\rg(f)=\dim\im f\leq\dim F=m$ car $\im f\subseteq F$, et
  $\rg(f)=n-\dim\Ker f\leq n$ car $\dim\Ker f\geq 0$.
\end{proof}

% ============================================================
\section{Injectivité, surjectivité, isomorphismes}\label{sec:inj-surj-iso}
% ============================================================

\begin{proposition}{Caractérisation de l'injectivité}{prop:inj-ker}
  Soit $f\colon E\to F$ une application linéaire. Alors
  \[
    f\text{ est injective}\iff\Ker f=\set{\vzero_E}.
  \]
\end{proposition}
\index{injectivité!application linéaire}

\begin{proof}
  $(\Rightarrow)$ Supposons $f$ injective. Soit $\vec{u}\in\Ker f$, alors
  $f(\vec{u})=\vzero_F=f(\vzero_E)$. Par injectivité, $\vec{u}=\vzero_E$.

  $(\Leftarrow)$ Supposons $\Ker f=\set{\vzero_E}$. Soient
  $\vec{u},\vec{v}\in E$ tels que $f(\vec{u})=f(\vec{v})$. Par linéarité,
  $f(\vec{u}-\vec{v})=\vzero_F$, donc $\vec{u}-\vec{v}\in\Ker f
  =\set{\vzero_E}$, d'où $\vec{u}=\vec{v}$.
\end{proof}

\begin{proposition}{Caractérisation de la surjectivité}{prop:surj-im}
  Soit $f\colon E\to F$ une application linéaire. Alors
  \[
    f\text{ est surjective}\iff\im f=F.
  \]
  En dimension finie, cela équivaut à $\rg(f)=\dim F$.
\end{proposition}
\index{surjectivité!application linéaire}

\begin{definition}{Isomorphisme}{def:isomorphisme}
  Un \emph{isomorphisme}\index{isomorphisme} d'espaces vectoriels est une
  application linéaire bijective. Si un isomorphisme $f\colon E\to F$
  existe, on dit que $E$ et $F$ sont \emph{isomorphes}\index{isomorphe} et
  l'on note $E\simeq F$ (ou $E\cong F$).
\end{definition}

\begin{theorem}{Bijectivité en dimension finie}{thm:bij-dim-finie}
  Soient $E$ et $F$ deux $\K$-espaces vectoriels de \textbf{même} dimension
  finie~$n$, et soit $f\colon E\to F$ linéaire. Les conditions suivantes
  sont équivalentes~:
  \begin{enumerate}[label=(\roman*)]
    \item $f$ est injective.
    \item $f$ est surjective.
    \item $f$ est bijective (isomorphisme).
  \end{enumerate}
\end{theorem}

\begin{proof}
  Posons $n=\dim E=\dim F$. Par le théorème du rang
  (
ef{thm:thm:rang}), $n=\dim\Ker f+\rg(f)$.

  (i)$\Rightarrow$(iii)~: Si $f$ est injective, $\dim\Ker f=0$, donc
  $\rg(f)=n=\dim F$, et $f$ est surjective, donc bijective.

  (ii)$\Rightarrow$(iii)~: Si $f$ est surjective, $\rg(f)=\dim F=n$,
  donc $\dim\Ker f=0$, et $f$ est injective, donc bijective.

  (iii)$\Rightarrow$(i) et (iii)$\Rightarrow$(ii) sont triviales.
\end{proof}

\begin{corollary}{Isomorphisme et dimension}{cor:iso-dim}
  Deux $\K$-espaces vectoriels de dimension finie sont isomorphes si et
  seulement s'ils ont la même dimension. En particulier, tout $\K$-espace
  vectoriel de dimension~$n$ est isomorphe à~$\K^n$.
\end{corollary}

\begin{proof}
  $(\Rightarrow)$ Si $f\colon E\isoto F$ est un isomorphisme, alors
  $\rg(f)=\dim F$ et $\dim\Ker f=0$, d'où $\dim E=\dim F$.

  $(\Leftarrow)$ Si $\dim E=\dim F=n$, choisissons une base
  $(\vece{1},\ldots,\vece{n})$ de~$E$ et une base
  $(\vecf{1},\ldots,\vecf{n})$ de~$F$. L'application $f$ définie par
  $f(\vece{i})=\vecf{i}$ pour tout~$i$ est un isomorphisme (exercice~: le
  vérifier).
\end{proof}

\begin{proposition}{L'inverse d'un isomorphisme est linéaire}{prop:inv-iso-lin}
  Si $f\colon E\to F$ est un isomorphisme, alors $\inv{f}\colon F\to E$
  est aussi une application linéaire (et donc un isomorphisme).
\end{proposition}

\begin{proof}
  Soient $\vec{w}_1,\vec{w}_2\in F$ et $\lambda\in\K$. Posons
  $\vec{u}_i=\inv{f}(\vec{w}_i)$, de sorte que $f(\vec{u}_i)=\vec{w}_i$.
  Alors
  \[
    f(\vec{u}_1+\lambda\vec{u}_2)
    =f(\vec{u}_1)+\lambda\,f(\vec{u}_2)
    =\vec{w}_1+\lambda\vec{w}_2.
  \]
  Donc $\inv{f}(\vec{w}_1+\lambda\vec{w}_2)
  =\vec{u}_1+\lambda\vec{u}_2
  =\inv{f}(\vec{w}_1)+\lambda\,\inv{f}(\vec{w}_2)$.
\end{proof}

\begin{definition}{Automorphisme}{def:automorphisme}
  Un \emph{automorphisme}\index{automorphisme} de $E$ est un isomorphisme
  de $E$ dans lui-même. L'ensemble des automorphismes de~$E$ est noté
  $\Aut(E)$ ou $\GL(E)$\index{GL(E)@$\GL(E)$}.
\end{definition}

% ============================================================
\section{Composition d'applications linéaires}\label{sec:composition-al}
% ============================================================

\begin{proposition}{La composée d'applications linéaires est linéaire}{prop:comp-lin}
  Soient $f\colon E\to F$ et $g\colon F\to G$ deux applications linéaires.
  Alors $g\circ f\colon E\to G$ est linéaire.
\end{proposition}

\begin{proof}
  Soient $\vec{u},\vec{v}\in E$ et $\lambda\in\K$. Alors
  \begin{align*}
    (g\circ f)(\vec{u}+\lambda\vec{v})
    &=g(f(\vec{u}+\lambda\vec{v}))
    =g(f(\vec{u})+\lambda\,f(\vec{v}))\\
    &=g(f(\vec{u}))+\lambda\,g(f(\vec{v}))
    =(g\circ f)(\vec{u})+\lambda\,(g\circ f)(\vec{v}).
  \end{align*}
\end{proof}

\begin{proposition}{Propriétés de la composition}{prop:comp-proprietes}
  Soient $f\colon E\to F$, $g\colon F\to G$ et $h\colon G\to H$ des
  applications linéaires.
  \begin{enumerate}[label=(\roman*)]
    \item \textbf{Associativité~:} $h\circ(g\circ f)=(h\circ g)\circ f$.
    \item \textbf{Élément neutre~:} $f\circ\Id_E=f$ et $\Id_F\circ f=f$.
    \item \textbf{Inégalité de Sylvester pour le rang~:}
      $\rg(g\circ f)\leq\min(\rg(f),\rg(g))$.
  \end{enumerate}
\end{proposition}

\begin{proof}
  Les points (i) et (ii) sont immédiats. Pour (iii), on a
  $\im(g\circ f)=g(f(E))\subseteq g(F)=\im g$, donc
  $\rg(g\circ f)\leq\rg(g)$. Par ailleurs, $\im(g\circ f)=g(\im f)$ et
  la restriction de~$g$ à $\im f$ a un rang au plus $\dim(\im f)=\rg(f)$,
  donc $\rg(g\circ f)\leq\rg(f)$.
\end{proof}

% ============================================================
\section{L'espace vectoriel $\cL(E,F)$}\label{sec:espace-LEF}
% ============================================================

\begin{theorem}{$\cL(E,F)$ est un espace vectoriel}{thm:LEF-ev}
  Soient $E$ et $F$ deux $\K$-espaces vectoriels. L'ensemble $\cL(E,F)$
  des applications linéaires de $E$ dans $F$, muni des opérations~:
  \begin{align*}
    (f+g)(\vec{u})&\deff f(\vec{u})+g(\vec{u}),\\
    (\lambda f)(\vec{u})&\deff\lambda\,f(\vec{u}),
  \end{align*}
  est un $\K$-espace vectoriel. L'élément nul est l'application nulle
  $\vec{u}\mapsto\vzero_F$.
\end{theorem}

\begin{proof}
  Il suffit de vérifier que $f+g$ et $\lambda f$ sont linéaires lorsque
  $f,g\in\cL(E,F)$ et $\lambda\in\K$, puis que les axiomes d'espace
  vectoriel sont satisfaits. La vérification est directe à partir de la
  structure d'espace vectoriel de~$F$.
\end{proof}

\begin{proposition}{Dimension de $\cL(E,F)$}{prop:dim-LEF}
  Si $\dim E=n$ et $\dim F=m$, alors
  \[
    \dim\cL(E,F)=nm.
  \]
\end{proposition}

\begin{proof}
  Choisissons des bases $(\vece{1},\ldots,\vece{n})$ de~$E$ et
  $(\vecf{1},\ldots,\vecf{m})$ de~$F$. Pour $1\leq i\leq n$ et
  $1\leq j\leq m$, définissons $f_{ij}\in\cL(E,F)$ par
  $f_{ij}(\vece{k})=\delta_{ik}\vecf{j}$ (où $\delta_{ik}$ est le
  symbole de Kronecker). La famille $(f_{ij})_{1\leq i\leq n,\,1\leq j\leq m}$
  forme une base de $\cL(E,F)$, donc $\dim\cL(E,F)=nm$.
\end{proof}

\begin{remark}{Structure de $\End(E)$}{rem:end-anneau}
  Lorsque $E=F$, l'ensemble $\End(E)=\cL(E,E)$ est non seulement un
  espace vectoriel, mais aussi un \emph{anneau} pour la composition~:
  la composition joue le rôle de la multiplication, $\Id_E$ est l'élément
  unité, et la distributivité vient de la linéarité. C'est un anneau non
  commutatif en général (dès que $\dim E\geq 2$).
\end{remark}

% ============================================================
\section{Matrice d'une application linéaire}\label{sec:matrice-al}
% ============================================================

\begin{definition}{Matrice associée à une application linéaire}{def:mat-al}
  Soient $E$ et $F$ deux $\K$-espaces vectoriels de dimensions finies
  $n$ et $m$ respectivement, munis de bases
  $\cB=(\vece{1},\ldots,\vece{n})$ et $\cC=(\vecf{1},\ldots,\vecf{m})$.
  Soit $f\in\cL(E,F)$.

  Pour chaque $j\in\set{1,\ldots,n}$, on décompose $f(\vece{j})$ dans la
  base~$\cC$~:
  \[
    f(\vece{j})=\sum_{i=1}^{m}a_{ij}\,\vecf{i}.
  \]
  La \emph{matrice de $f$ relativement aux bases $\cB$ et $\cC$}%
  \index{matrice!d'une application linéaire} est la matrice
  $A=(a_{ij})\in\Mat_{m,n}(\K)$ dont la $j$-ème colonne contient les
  coordonnées de $f(\vece{j})$ dans la base~$\cC$. On la note
  \[
    \Mat_{\cB,\cC}(f)\deff A.
  \]
  Lorsque $E=F$ et $\cB=\cC$, on écrit simplement $\Mat_\cB(f)$.
\end{definition}

\begin{proposition}{Lien entre application linéaire et produit matriciel}{prop:al-prod-mat}
  Avec les notations précédentes, si $\vec{u}=\sum_{j=1}^n x_j\,\vece{j}$
  a pour vecteur de coordonnées $X=(x_1,\ldots,x_n)^{\mathsf{T}}\in\K^n$
  dans la base~$\cB$, alors $f(\vec{u})$ a pour vecteur de coordonnées
  $Y=AX\in\K^m$ dans la base~$\cC$, où $A=\Mat_{\cB,\cC}(f)$.
\end{proposition}

\begin{proof}
  On a $f(\vec{u})=f\!\left(\sum_{j=1}^n x_j\,\vece{j}\right)
  =\sum_{j=1}^n x_j\,f(\vece{j})
  =\sum_{j=1}^n x_j\sum_{i=1}^m a_{ij}\,\vecf{i}
  =\sum_{i=1}^m\left(\sum_{j=1}^n a_{ij}x_j\right)\vecf{i}$.
  La $i$-ème coordonnée de $f(\vec{u})$ dans $\cC$ est donc
  $y_i=\sum_{j=1}^n a_{ij}x_j$, ce qui est exactement la $i$-ème
  composante du produit $AX$.
\end{proof}

\begin{example}{Matrice de la rotation}{ex:mat-rotation}
  La rotation $R_\theta$ de $\R^2$ dans la base canonique
  $\cE=(\vece{1},\vece{2})$ a pour matrice~:
  \[
    \Mat_\cE(R_\theta)=\begin{pmatrix}
      \cos\theta & -\sin\theta \\ \sin\theta & \cos\theta
    \end{pmatrix}.
  \]
  En effet, $R_\theta(\vece{1})=\begin{pmatrix}\cos\theta\\\sin\theta\end{pmatrix}$
  et $R_\theta(\vece{2})=\begin{pmatrix}-\sin\theta\\\cos\theta\end{pmatrix}$.
\end{example}

\begin{example}{Matrice de la projection}{ex:mat-projection}
  La projection orthogonale $p$ sur l'axe des abscisses dans $\R^2$ a
  pour matrice dans la base canonique~:
  \[
    \Mat_\cE(p)=\begin{pmatrix}1&0\\0&0\end{pmatrix}.
  \]
\end{example}

\begin{example}{Matrice de la dérivation}{ex:mat-derivation}
  Soit $D\colon\R_3[X]\to\R_3[X]$ la dérivation, dans la base
  $\cB=(1,X,X^2,X^3)$. On a $D(1)=0$, $D(X)=1$, $D(X^2)=2X$,
  $D(X^3)=3X^2$. Donc~:
  \[
    \Mat_\cB(D)=\begin{pmatrix}
      0&1&0&0\\ 0&0&2&0\\ 0&0&0&3\\ 0&0&0&0
    \end{pmatrix}.
  \]
\end{example}

\begin{theorem}{Isomorphisme $\cL(E,F)\simeq\Mat_{m,n}(\K)$}{thm:iso-LEF-mat}
  L'application
  \[
    \Phi\colon\cL(E,F)\to\Mat_{m,n}(\K),\qquad f\mapsto\Mat_{\cB,\cC}(f)
  \]
  est un isomorphisme d'espaces vectoriels. De plus~:
  \begin{enumerate}[label=(\roman*)]
    \item $\Mat_{\cB,\cC}(f+g)=\Mat_{\cB,\cC}(f)+\Mat_{\cB,\cC}(g)$.
    \item $\Mat_{\cB,\cC}(\lambda f)=\lambda\,\Mat_{\cB,\cC}(f)$.
    \item Si $g\colon F\to G$ et $\cD$ est une base de~$G$, alors
      $\Mat_{\cB,\cD}(g\circ f)=\Mat_{\cC,\cD}(g)\cdot\Mat_{\cB,\cC}(f)$.
  \end{enumerate}
\end{theorem}

\begin{proof}
  La linéarité de $\Phi$ découle des points (i) et (ii), qui se
  vérifient par identification des coefficients. L'injectivité vient du
  fait que si $\Mat_{\cB,\cC}(f)=0$, alors $f(\vece{j})=\vzero$ pour
  tout~$j$, donc $f=0$. La surjectivité vient du fait que toute matrice
  $A$ définit une application linéaire $f$ par
  $f(\vece{j})=\sum_i a_{ij}\vecf{i}$.

  Pour (iii), les colonnes de $\Mat_{\cB,\cD}(g\circ f)$ sont les
  coordonnées de $(g\circ f)(\vece{j})=g(f(\vece{j}))$ dans~$\cD$.
  Si $\Mat_{\cB,\cC}(f)=A$ et $\Mat_{\cC,\cD}(g)=B$, alors
  $f(\vece{j})=\sum_k a_{kj}\vecf{k}$ et
  $g(f(\vece{j}))=\sum_k a_{kj}g(\vecf{k})
  =\sum_k a_{kj}\sum_i b_{ik}\vec{d}_i
  =\sum_i\left(\sum_k b_{ik}a_{kj}\right)\vec{d}_i$,
  ce qui donne la matrice $BA$.
\end{proof}

% ============================================================
\section{Changement de base}\label{sec:changement-base}
% ============================================================

\begin{definition}{Matrice de passage}{def:mat-passage}
  Soient $\cB=(\vece{1},\ldots,\vece{n})$ et
  $\cB'=(\vece{1}',\ldots,\vece{n}')$ deux bases de~$E$. La
  \emph{matrice de passage}\index{matrice de passage} de $\cB$ à $\cB'$
  est la matrice
  \[
    P_{\cB\to\cB'}\deff\Mat_\cB(\Id_E)_{\cB,\cB'}
  \]
  dont la $j$-ème colonne contient les coordonnées de $\vece{j}'$ dans la
  base~$\cB$. Autrement dit, si
  $\vece{j}'=\sum_{i=1}^n p_{ij}\,\vece{i}$, alors $P=(p_{ij})$.
\end{definition}

\begin{proposition}{Propriétés des matrices de passage}{prop:mat-passage}
  \begin{enumerate}[label=(\roman*)]
    \item $P_{\cB\to\cB'}$ est inversible et
      $\inv{(P_{\cB\to\cB'})}=P_{\cB'\to\cB}$.
    \item Si $X$ est le vecteur de coordonnées de $\vec{u}$ dans $\cB$ et
      $X'$ est celui dans $\cB'$, alors $X=P_{\cB\to\cB'}\,X'$, soit
      $X'=\inv{(P_{\cB\to\cB'})}\,X$.
    \item Pour trois bases $\cB,\cB',\cB''$~:
      $P_{\cB\to\cB''}=P_{\cB\to\cB'}\cdot P_{\cB'\to\cB''}$.
  \end{enumerate}
\end{proposition}

\begin{proof}
  \emph{(i)} La matrice $P_{\cB\to\cB'}$ est la matrice de l'identité
  $\Id_E$ de la base $\cB'$ vers la base $\cB$. Comme $\Id_E$ est un
  isomorphisme, sa matrice est inversible. De plus,
  $P_{\cB\to\cB'}\cdot P_{\cB'\to\cB}$ est la matrice de
  $\Id_E\circ\Id_E=\Id_E$ dans la base $\cB$, qui est la matrice
  identité~$I_n$.

  \emph{(ii)} Soit $\vec{u}=\sum_i x_i\vece{i}=\sum_j x_j'\vece{j}'$.
  En remplaçant $\vece{j}'=\sum_i p_{ij}\vece{i}$, on obtient
  $\vec{u}=\sum_i\left(\sum_j p_{ij}x_j'\right)\vece{i}$, d'où
  $x_i=\sum_j p_{ij}x_j'$, c'est-à-dire $X=PX'$.

  \emph{(iii)} Résulte de la multiplicativité du passage~:
  $X=P_{\cB\to\cB'}\,X'$ et $X'=P_{\cB'\to\cB''}\,X''$, d'où
  $X=P_{\cB\to\cB'}\cdot P_{\cB'\to\cB''}\,X''$.
\end{proof}

\begin{theorem}{Formule de changement de base}{thm:changement-base}
  Soit $f\in\End(E)$. Soient $\cB$ et $\cB'$ deux bases de~$E$, et soit
  $P=P_{\cB\to\cB'}$ la matrice de passage. Si $A=\Mat_\cB(f)$ est la
  matrice de $f$ dans $\cB$, alors la matrice de $f$ dans $\cB'$ est
  \[
    \boxed{A'=\Mat_{\cB'}(f)=\inv{P}\,A\,P.}
  \]
\end{theorem}
\index{changement de base}

\begin{proof}
  La matrice de $f$ dans $\cB'$ est $\Mat_{\cB'}(f)
  =\Mat_{\cB',\cB'}(\Id_E\circ f\circ\Id_E)$. Or $\Id_E$ de $\cB'$
  vers $\cB$ a pour matrice $P$ et de $\cB$ vers $\cB'$ a pour matrice
  $\inv{P}$. Par la formule de composition des matrices
  (
ef{thm:thm:iso-LEF-mat}, (iii))~:
  \[
    \Mat_{\cB'}(f)=\inv{P}\cdot\Mat_\cB(f)\cdot P=\inv{P}AP. \qedhere
  \]
\end{proof}

\begin{definition}{Matrices semblables}{def:semblables}
  Deux matrices $A,B\in\Mat_n(\K)$ sont dites
  \emph{semblables}\index{matrices semblables} s'il existe une matrice
  inversible $P\in\GL_n(\K)$ telle que $B=\inv{P}AP$. Les matrices
  semblables représentent le même endomorphisme dans des bases
  différentes.
\end{definition}

\begin{remark}{Cas général : deux espaces différents}{rem:chgt-base-general}
  Si $f\colon E\to F$ avec $\cB,\cB'$ bases de~$E$ et $\cC,\cC'$ bases
  de~$F$, et si $P=P_{\cB\to\cB'}$, $Q=P_{\cC\to\cC'}$, alors~:
  \[
    \Mat_{\cB',\cC'}(f)=\inv{Q}\cdot\Mat_{\cB,\cC}(f)\cdot P.
  \]
\end{remark}

% ============================================================
\section{Opérations sur les matrices}\label{sec:operations-matrices}
% ============================================================

Dans cette section, on rappelle les opérations algébriques fondamentales
sur les matrices et leurs propriétés.

\subsection{Addition et multiplication scalaire}\label{subsec:add-scal-mat}

\begin{definition}{Addition et multiplication scalaire}{def:add-scal-mat}
  Soient $A=(a_{ij}),B=(b_{ij})\in\Mat_{m,n}(\K)$ et $\lambda\in\K$.
  \begin{enumerate}[label=(\roman*)]
    \item \textbf{Addition~:} $A+B\deff(a_{ij}+b_{ij})$.
    \item \textbf{Multiplication scalaire~:} $\lambda A\deff(\lambda\,a_{ij})$.
  \end{enumerate}
  L'ensemble $\Mat_{m,n}(\K)$ muni de ces opérations est un $\K$-espace
  vectoriel de dimension~$mn$.
\end{definition}

\subsection{Produit matriciel}\label{subsec:prod-mat}

\begin{definition}{Produit de deux matrices}{def:prod-mat}
  Soient $A=(a_{ij})\in\Mat_{m,p}(\K)$ et $B=(b_{jk})\in\Mat_{p,n}(\K)$.
  Le \emph{produit}\index{produit matriciel} $AB$ est la matrice
  $C=(c_{ik})\in\Mat_{m,n}(\K)$ définie par
  \[
    c_{ik}=\sum_{j=1}^{p}a_{ij}b_{jk},\qquad
    1\leq i\leq m,\;1\leq k\leq n.
  \]
\end{definition}

\begin{proposition}{Propriétés du produit matriciel}{prop:prod-mat}
  Pour des matrices de tailles compatibles~:
  \begin{enumerate}[label=(\roman*)]
    \item \textbf{Associativité~:} $(AB)C=A(BC)$.
    \item \textbf{Distributivité~:} $A(B+C)=AB+AC$ et $(A+B)C=AC+BC$.
    \item \textbf{Élément neutre~:} $I_n A=A$ et $AI_m=A$ (pour
      $A\in\Mat_{n,m}(\K)$), où $I_k$ est la matrice identité.
    \item \textbf{Compatibilité scalaire~:} $\lambda(AB)=(\lambda A)B=A(\lambda B)$.
    \item \textbf{Non-commutativité en général~:} $AB\neq BA$ en général.
  \end{enumerate}
\end{proposition}

\begin{example}{Non-commutativité}{ex:non-comm}
  Soient $A=\begin{pmatrix}1&1\\0&0\end{pmatrix}$ et
  $B=\begin{pmatrix}0&0\\1&1\end{pmatrix}$. Alors
  \[
    AB=\begin{pmatrix}1&1\\0&0\end{pmatrix},\qquad
    BA=\begin{pmatrix}0&0\\1&1\end{pmatrix}.
  \]
  On a $AB\neq BA$.
\end{example}

\subsection{Transposition}\label{subsec:transposition}

\begin{definition}{Transposée}{def:transposee}
  Soit $A=(a_{ij})\in\Mat_{m,n}(\K)$. La \emph{transposée}%
  \index{transposée} de~$A$ est la matrice
  $\transpose{A}=(a_{ji})\in\Mat_{n,m}(\K)$ obtenue en échangeant lignes
  et colonnes.
\end{definition}

\begin{proposition}{Propriétés de la transposée}{prop:transposee}
  Pour $A,B$ de tailles compatibles et $\lambda\in\K$~:
  \begin{enumerate}[label=(\roman*)]
    \item $\transpose{(A+B)}=\transpose{A}+\transpose{B}$.
    \item $\transpose{(\lambda A)}=\lambda\,\transpose{A}$.
    \item $\transpose{(\transpose{A})}=A$.
    \item $\transpose{(AB)}=\transpose{B}\cdot\transpose{A}$
      (inversion de l'ordre).
  \end{enumerate}
\end{proposition}

\begin{proof}
  Les points (i)--(iii) sont immédiats par définition. Pour (iv), le
  coefficient $(i,j)$ de $\transpose{(AB)}$ est $(AB)_{ji}
  =\sum_k a_{jk}b_{ki}=\sum_k (\transpose{B})_{ik}(\transpose{A})_{kj}
  =(\transpose{B}\cdot\transpose{A})_{ij}$.
\end{proof}

% ============================================================
\section{Matrices inversibles}\label{sec:matrices-inversibles}
% ============================================================

\begin{definition}{Matrice inversible}{def:mat-inv}
  Une matrice $A\in\Mat_n(\K)$ est dite
  \emph{inversible}\index{matrice inversible} s'il existe une matrice
  $B\in\Mat_n(\K)$ telle que $AB=BA=I_n$. On note alors $B=\inv{A}$.
  L'ensemble des matrices inversibles d'ordre~$n$ est noté
  $\GL_n(\K)$\index{GL@$\GL_n(\K)$}.
\end{definition}

\begin{proposition}{Caractérisations de l'inversibilité}{prop:carac-inv}
  Pour $A\in\Mat_n(\K)$, les conditions suivantes sont équivalentes~:
  \begin{enumerate}[label=(\roman*)]
    \item $A$ est inversible.
    \item L'application $\vec{x}\mapsto A\vec{x}$ de $\K^n$ dans $\K^n$
      est bijective.
    \item Les colonnes de $A$ forment une base de~$\K^n$.
    \item $\rg(A)=n$.
    \item $\Ker(A)=\set{\vzero}$ (où l'on identifie $A$ à l'application
      linéaire correspondante).
  \end{enumerate}
\end{proposition}

\begin{proposition}{Propriétés de l'inverse}{prop:proprietes-inv}
  Soient $A,B\in\GL_n(\K)$ et $\lambda\in\K^*$.
  \begin{enumerate}[label=(\roman*)]
    \item $\inv{(AB)}=\inv{B}\cdot\inv{A}$.
    \item $\inv{(\inv{A})}=A$.
    \item $\inv{(\lambda A)}=\lambda^{-1}\inv{A}$.
    \item $\inv{(\transpose{A})}=\transpose{(\inv{A})}$.
  \end{enumerate}
\end{proposition}

\begin{proof}
  \emph{(i)} $(AB)(\inv{B}\inv{A})=A(B\inv{B})\inv{A}=A\,I_n\,\inv{A}
  =A\inv{A}=I_n$. De même pour le produit dans l'autre sens.

  \emph{(ii)} Résulte de $\inv{A}\cdot A=I_n$ et $A\cdot\inv{A}=I_n$.

  \emph{(iii)} $(\lambda A)(\lambda^{-1}\inv{A})
  =(\lambda\cdot\lambda^{-1})(A\inv{A})=I_n$.

  \emph{(iv)} $\transpose{A}\cdot\transpose{(\inv{A})}
  =\transpose{(\inv{A}\cdot A)}=\transpose{I_n}=I_n$.
\end{proof}

\begin{example}{Inverse d'une matrice $2\times 2$}{ex:inv-2x2}
  Soit $A=\begin{pmatrix}a&b\\c&d\end{pmatrix}$ avec $ad-bc\neq 0$.
  Alors
  \[
    \inv{A}=\frac{1}{ad-bc}\begin{pmatrix}d&-b\\-c&a\end{pmatrix}.
  \]
  La quantité $\det A=ad-bc$ est le \emph{déterminant}\index{déterminant}
  de~$A$ (nous l'étudierons en détail au chapitre suivant).
\end{example}

\begin{remark}{Le groupe $\GL_n(\K)$}{rem:GLn}
  L'ensemble $\GL_n(\K)$, muni de la multiplication matricielle, est un
  groupe (non abélien pour $n\geq 2$). C'est le \emph{groupe linéaire
  général}\index{groupe linéaire}.
\end{remark}

% ============================================================
\section{Espaces duaux (introduction)}\label{sec:dual}
% ============================================================

\begin{definition}{Espace dual}{def:dual}
  Soit $E$ un $\K$-espace vectoriel. L'\emph{espace dual}%
  \index{espace dual}\index{dual} de~$E$ est l'espace
  \[
    E^*\deff\cL(E,\K).
  \]
  Les éléments de $E^*$ sont appelés \emph{formes linéaires}%
  \index{forme linéaire} sur~$E$.
\end{definition}

\begin{example}{Formes linéaires sur $\R^n$}{ex:forme-lin-Rn}
  Toute forme linéaire sur $\R^n$ est de la forme
  $\vphi(\vec{x})=a_1 x_1+a_2 x_2+\cdots+a_n x_n$ pour des scalaires
  $a_1,\ldots,a_n\in\R$ fixés. On peut l'écrire
  $\vphi(\vec{x})=\transpose{\vec{a}}\,\vec{x}$ où
  $\vec{a}=(a_1,\ldots,a_n)^{\mathsf{T}}$.
\end{example}

\begin{definition}{Base duale}{def:base-duale}
  Soit $\cB=(\vece{1},\ldots,\vece{n})$ une base de~$E$ (de dimension
  finie). La \emph{base duale}\index{base duale} de~$\cB$ est la famille
  $\cB^*=(\vece{1}^*,\ldots,\vece{n}^*)$ de formes linéaires définies
  par
  \[
    \vece{i}^*(\vece{j})=\delta_{ij}=
    \begin{cases}1&\text{si }i=j,\\0&\text{si }i\neq j.\end{cases}
  \]
\end{definition}

\begin{proposition}{La base duale est une base de $E^*$}{prop:base-duale}
  La famille $\cB^*=(\vece{1}^*,\ldots,\vece{n}^*)$ est une base de
  $E^*$. En particulier, $\dim E^*=\dim E$, et donc $E^*\simeq E$.

  Toute forme linéaire $\vphi\in E^*$ se décompose de manière unique~:
  \[
    \vphi=\sum_{i=1}^n\vphi(\vece{i})\,\vece{i}^*.
  \]
\end{proposition}

\begin{proof}
  \emph{Liberté.} Supposons $\sum_{i=1}^n\lambda_i\,\vece{i}^*=0$ (la
  forme nulle). En évaluant en $\vece{j}$~:
  $0=\sum_{i=1}^n\lambda_i\,\vece{i}^*(\vece{j})
  =\sum_{i=1}^n\lambda_i\delta_{ij}=\lambda_j$. Donc $\lambda_j=0$
  pour tout~$j$.

  \emph{Famille génératrice.} Soit $\vphi\in E^*$ et posons
  $\psi=\sum_{i=1}^n\vphi(\vece{i})\,\vece{i}^*$. Pour tout~$j$~:
  $\psi(\vece{j})=\sum_{i=1}^n\vphi(\vece{i})\,\delta_{ij}
  =\vphi(\vece{j})$. Comme $\vphi$ et $\psi$ coïncident sur une base
  de~$E$, elles sont égales.
\end{proof}

\begin{remark}{Bidual}{rem:bidual}
  Le \emph{bidual}\index{bidual} $E^{**}=(E^*)^*$ contient les formes
  linéaires sur les formes linéaires. En dimension finie, l'application
  canonique $\iota\colon E\to E^{**}$, définie par
  $\iota(\vec{u})(\vphi)=\vphi(\vec{u})$, est un isomorphisme naturel
  (qui ne dépend pas du choix d'une base). On identifie souvent $E$ à
  $E^{**}$ par cet isomorphisme.
\end{remark}

% ============================================================
\section{Applications}\label{sec:applications-ch2}
% ============================================================

\subsection{Transformations en infographie}\label{subsec:infographie}

En infographie 2D et 3D\index{infographie}, les transformations
géométriques (translations, rotations, homothéties, projections) sont
réalisées par multiplication matricielle. Pour pouvoir exprimer les
translations (qui ne sont pas linéaires) sous forme matricielle, on
utilise les \emph{coordonnées homogènes}\index{coordonnées homogènes}~:
un point $(x,y)$ du plan est représenté par le vecteur
$(x,y,1)^{\mathsf{T}}\in\R^3$.

\begin{example}{Transformations 2D en coordonnées homogènes}{ex:transfo-2d}
  \begin{enumerate}[label=(\alph*)]
    \item \textbf{Translation} de vecteur $(t_x,t_y)$~:
    $T=\begin{pmatrix}1&0&t_x\\0&1&t_y\\0&0&1\end{pmatrix}$.
    \item \textbf{Rotation} d'angle $\theta$~:
    $R=\begin{pmatrix}\cos\theta&-\sin\theta&0\\\sin\theta&\cos\theta&0\\0&0&1\end{pmatrix}$.
    \item \textbf{Homothétie} de facteurs $(s_x,s_y)$~:
    $S=\begin{pmatrix}s_x&0&0\\0&s_y&0\\0&0&1\end{pmatrix}$.
  \end{enumerate}
  La composition de ces transformations revient à multiplier les matrices.
  Dans un moteur graphique, l'ensemble des transformations appliquées à un
  objet est stocké dans une seule matrice.
\end{example}

\subsection{Aperçu des chaînes de Markov}\label{subsec:markov}

\begin{example}{Chaîne de Markov simple}{ex:markov}
  Un système peut se trouver dans deux états $S_1$ ou $S_2$. À chaque
  étape, la probabilité de transition de $S_i$ à $S_j$ est donnée par la
  \emph{matrice de transition}\index{matrice de transition}~:
  \[
    P=\begin{pmatrix}0{,}7&0{,}4\\0{,}3&0{,}6\end{pmatrix},
  \]
  où $p_{ij}$ est la probabilité de passer de l'état $j$ à l'état $i$.
  Chaque colonne somme à~$1$ (matrice \emph{stochastique}\index{matrice
  stochastique}). Si $\vec{x}_k$ est le vecteur de probabilités d'état à
  l'étape~$k$, alors $\vec{x}_{k+1}=P\vec{x}_k$. Les puissances $P^k$
  décrivent l'évolution à long terme. La recherche d'un \emph{état
  stationnaire} $\vec{x}$ tel que $P\vec{x}=\vec{x}$ revient à chercher
  un vecteur propre associé à la valeur propre~$1$ --- sujet que nous
  approfondirons dans un chapitre ultérieur.
\end{example}

% ============================================================
\section{Exercices}\label{sec:exo-ch2}
% ============================================================

\begin{exercise}{\basic{} Vérification de linéarité}{exo:verif-lin}
  Pour chacune des applications suivantes, dire si elle est linéaire.
  Justifier.
  \begin{enumerate}[label=(\alph*)]
    \item $f\colon\R^2\to\R$, $(x,y)\mapsto 2x-3y$.
    \item $g\colon\R^2\to\R$, $(x,y)\mapsto x^2+y$.
    \item $h\colon\R^2\to\R^2$, $(x,y)\mapsto(x+1,\,y)$.
    \item $\vphi\colon\R^3\to\R^2$, $(x,y,z)\mapsto(x-z,\,2y+z)$.
    \item $\psi\colon\R_2[X]\to\R_2[X]$, $P\mapsto P-P(0)$.
  \end{enumerate}
\end{exercise}

\begin{exercise}{\basic{} Noyau et image}{exo:ker-im}
  Soit $f\colon\R^3\to\R^2$ définie par
  $f(x,y,z)=(x+y-z,\;2x+y+z)$.
  \begin{enumerate}[label=(\alph*)]
    \item Montrer que $f$ est linéaire.
    \item Déterminer $\Ker f$ (en donner une base).
    \item Déterminer $\im f$.
    \item Vérifier le théorème du rang.
  \end{enumerate}
\end{exercise}

\begin{exercise}{\basic{} Matrice d'une application linéaire}{exo:mat-al}
  Soit $f\colon\R^3\to\R^2$ définie par $f(x,y,z)=(x+2y,\;3y-z)$.
  \begin{enumerate}[label=(\alph*)]
    \item Écrire la matrice de $f$ dans les bases canoniques.
    \item Calculer $f(1,-1,2)$ de deux manières~: directement et par
      produit matriciel.
    \item Déterminer le rang de $f$.
  \end{enumerate}
\end{exercise}

\begin{exercise}{\basic{} Projections}{exo:projections}
  Soit $p\colon\R^3\to\R^3$ définie par $p(x,y,z)=(x,y,0)$.
  \begin{enumerate}[label=(\alph*)]
    \item Montrer que $p$ est linéaire.
    \item Vérifier que $p\circ p=p$.
    \item Déterminer $\Ker p$ et $\im p$.
    \item Montrer que $\R^3=\Ker p\oplus\im p$.
  \end{enumerate}
\end{exercise}

\begin{exercise}{\intermediate{} Injectivité et surjectivité}{exo:inj-surj-al}
  Soit $f\colon\R^3\to\R^3$ l'endomorphisme défini par
  $f(x,y,z)=(x+y,\;y+z,\;x+z)$.
  \begin{enumerate}[label=(\alph*)]
    \item Écrire la matrice de $f$ dans la base canonique.
    \item $f$ est-elle injective~? Surjective~? Bijective~?
    \item Si $f$ est bijective, trouver $\inv{f}$.
  \end{enumerate}
\end{exercise}

\begin{exercise}{\intermediate{} Dérivation et intégration}{exo:deriv-integ}
  Soit $E=\R_3[X]$ muni de la base $\cB=(1,X,X^2,X^3)$.
  \begin{enumerate}[label=(\alph*)]
    \item Écrire la matrice de $D\colon P\mapsto P'$ dans $\cB$.
    \item Écrire la matrice de $I\colon P\mapsto\int_0^X P(t)\,dt$ dans
      $\cB$ (l'image est dans $\R_4[X]$~; prendre la base
      $\cC=(1,X,X^2,X^3,X^4)$).
    \item Calculer le produit $DI$ (comme endomorphisme de $E$ en
      restreignant au degré $\leq 3$). Que constate-t-on~?
  \end{enumerate}
\end{exercise}

\begin{exercise}{\intermediate{} Changement de base}{exo:chgt-base}
  Soit $f\colon\R^2\to\R^2$ l'endomorphisme de matrice
  $A=\begin{pmatrix}3&1\\1&3\end{pmatrix}$ dans la base canonique $\cE$.
  Soit $\cB'=\left(\begin{pmatrix}1\\1\end{pmatrix},
  \begin{pmatrix}1\\-1\end{pmatrix}\right)$.
  \begin{enumerate}[label=(\alph*)]
    \item Écrire la matrice de passage $P$ de $\cE$ à $\cB'$.
    \item Calculer $\inv{P}$.
    \item Calculer $A'=\inv{P}AP$. Que remarque-t-on~?
    \item Interpréter géométriquement l'action de $f$.
  \end{enumerate}
\end{exercise}

\begin{exercise}{\intermediate{} Espace dual}{exo:dual}
  Soit $E=\R^3$ muni de la base canonique $\cE=(\vece{1},\vece{2},\vece{3})$.
  \begin{enumerate}[label=(\alph*)]
    \item Écrire la base duale $\cE^*=(\vece{1}^*,\vece{2}^*,\vece{3}^*)$
      explicitement (comme formes linéaires).
    \item Soit $\vphi\colon\R^3\to\R$, $(x,y,z)\mapsto 2x-y+3z$.
      Exprimer $\vphi$ dans la base duale.
    \item Soit $\cB=((1,1,0),(1,0,1),(0,1,1))$. Calculer la base duale
      $\cB^*$.
  \end{enumerate}
\end{exercise}

\begin{exercise}{\challenging{} Théorème du rang --- applications}{exo:thm-rang-app}
  \begin{enumerate}[label=(\alph*)]
    \item Soit $f\in\End(E)$ avec $\dim E=n$. Montrer que
      $\Ker f\cap\im f=\set{\vzero}$ si et seulement si
      $\Ker f\oplus\im f=E$.
    \item Soit $f\in\End(E)$ tel que $f\circ f=f$ (projecteur). Montrer
      que $E=\Ker f\oplus\im f$.
    \item Soit $f\in\End(E)$ tel que $\rg(f)=1$. Montrer qu'il existe
      $\lambda\in\K$ tel que $f^2=\lambda f$.

      \emph{Indication~:} choisir une base adaptée au noyau et à l'image.
  \end{enumerate}
\end{exercise}

\begin{exercise}{\challenging{} Matrices et applications linéaires}{exo:mat-avance}
  Soient $A\in\Mat_{m,n}(\K)$ et $B\in\Mat_{n,p}(\K)$.
  \begin{enumerate}[label=(\alph*)]
    \item Montrer que $\rg(AB)\leq\min(\rg(A),\rg(B))$.
    \item Montrer que $\rg(A)+\rg(B)-n\leq\rg(AB)$ (\emph{inégalité de
      Sylvester}\index{inégalité de Sylvester}).

      \emph{Indication~:} appliquer le théorème du rang à la restriction
      de $f_A$ à $\im(f_B)$.
    \item En déduire que si $A\in\Mat_n(\K)$ avec $\rg(A)=n-1$, alors
      $\rg(A^2)\geq n-2$.
  \end{enumerate}
\end{exercise}

\begin{exercise}{\challenging{} Endomorphismes nilpotents}{exo:nilpotent}
  Un endomorphisme $f\in\End(E)$ est dit
  \emph{nilpotent}\index{nilpotent} s'il existe $k\in\N^*$ tel que
  $f^k=0$ (l'application nulle). Le plus petit tel~$k$ est l'\emph{indice
  de nilpotence}.
  \begin{enumerate}[label=(\alph*)]
    \item Montrer que si $f$ est nilpotent d'indice~$k$, alors
      $k\leq\dim E$.

      \emph{Indication~:} considérer la suite
      $\set{\vzero}\subseteq\Ker f\subseteq\Ker f^2\subseteq\cdots$
    \item Montrer que si $f$ est nilpotent, alors $\Id+f$ est un
      automorphisme.

      \emph{Indication~:} écrire l'inverse comme une somme finie
      $\sum_{i=0}^{k-1}(-1)^i f^i$.
    \item Soit $E=\R^3$ et $f$ défini par $f(x,y,z)=(y,z,0)$. Vérifier
      que $f$ est nilpotent, déterminer son indice, et calculer
      $\inv{(\Id+f)}$.
  \end{enumerate}
\end{exercise}

% ============================================================
\section*{Résumé du chapitre}
\addcontentsline{toc}{section}{Résumé du chapitre}
% ============================================================

\begin{framed}
\textbf{Ce qu'il faut retenir.}

\begin{enumerate}[label=\textbf{\arabic*.},leftmargin=2em]
  \item \textbf{Application linéaire.} $f\colon E\to F$ est linéaire si
    $f(\lambda\vec{u}+\mu\vec{v})=\lambda f(\vec{u})+\mu f(\vec{v})$.
    Elle est entièrement déterminée par les images d'une base.

  \item \textbf{Noyau et image.} $\Ker f$ et $\im f$ sont des
    sous-espaces vectoriels. $f$ est injective ssi $\Ker f=\set{\vzero}$.

  \item \textbf{Théorème du rang.}
    $\dim E=\dim\Ker f+\rg(f)$.
    C'est l'outil principal pour déterminer injectivité et surjectivité
    en dimension finie.

  \item \textbf{Isomorphismes.} En dimension finie égale, injectivité
    $\Leftrightarrow$ surjectivité $\Leftrightarrow$ bijectivité. Deux
    espaces de même dimension sont isomorphes.

  \item \textbf{Matrice d'une application linéaire.} La $j$-ème colonne
    de $\Mat_{\cB,\cC}(f)$ contient les coordonnées de $f(\vece{j})$ dans
    $\cC$. La composition correspond au produit matriciel.

  \item \textbf{Changement de base.} Si $P$ est la matrice de passage,
    $A'=\inv{P}AP$. Deux matrices sont semblables ssi elles représentent
    le même endomorphisme.

  \item \textbf{Matrices inversibles.} $A\in\GL_n(\K)$ ssi
    $\rg(A)=n$ ssi $\Ker A=\set{\vzero}$ ssi les colonnes de $A$ forment
    une base.

  \item \textbf{Espace dual.} $E^*=\cL(E,\K)$ est l'espace des formes
    linéaires. En dimension finie, $\dim E^*=\dim E$.

  \item \textbf{Formules clés.}
    \begin{itemize}
      \item $\dim\cL(E,F)=\dim E\cdot\dim F$.
      \item $\rg(g\circ f)\leq\min(\rg f,\rg g)$.
      \item $\inv{(AB)}=\inv{B}\cdot\inv{A}$, \;
        $\transpose{(AB)}=\transpose{B}\cdot\transpose{A}$.
    \end{itemize}
\end{enumerate}
\end{framed}
